\chapter{Prime GNS Normalization and Arithmetic Pairings}
\label{v4-arith:kms-gns-full}

The finite primary character determines a partition trace. A
Koszul--Petersson comparison additionally requires a Hilbert
representation of the relevant observables and a map from the
native pairing to its state. We first construct the local GNS
representation and its modular conjugation with the normalization
of \Cref{v4-arith:thermal:chapter}. We then identify the further
comparison data needed for the arithmetic pairing problem.
The operator calculation is valid at every real $\beta>0$;
$\beta=1$ is a regular single-prime temperature.

\section{The Local Comparison Problem}
\label{v4-arith:kms-gns:statement-table}
\label{v4-arith:kms-gns:residuals}
\label{v4-arith:kms-gns:consequences}
\label{v4-arith:kms-gns:P3-consequence}
\label{v4-arith:kms-gns:table}

Three constructions are required by the original arithmetic
problem. The first is the state and its GNS modular involution.
The second is a pairing-preserving map from the native local
Koszul and Whittaker data to that representation, including
ramified factors. The third is an adelic assembly compatible
with the genuine arithmetic phases and the continuous-spectrum
domains. The first construction below does not supply the other
two. Temperedness and a local functional equation alone do not
construct a map from a categorical Hochschild trace to a Hilbert
space of arithmetic observables.

\section{The Local State and Its Represented Quotient}
\label{v4-arith:kms-gns:cycle1}
\label{v4-arith:kms-gns:cycle1:state}

Let $\Gamma_p=\mathbb Z[1/p]/\mathbb Z$. Its natural map to
$\mathbb Q_p/\mathbb Z_p$ is an isomorphism: the finitely many
negative $p$-adic digits give a rational representative with
$p$-power denominator, and such a rational is $p$-integral
precisely when its class in $\Gamma_p$ is zero. Consider the unital
complex $*$-algebra $\mathcal B_p$ presented by an isometry $u$
and unitaries $e(\gamma)$, $\gamma\in\Gamma_p$, with
\[
 u^*u=1,\quad e(\gamma)e(\eta)=e(\gamma+\eta),\quad
 e(0)=1,\quad e(\gamma)^*=e(-\gamma),
\]
\[
 u^*e(\gamma)u=e(p\gamma),\qquad
 ue(\gamma)u^*=\frac1p\sum_{p\eta=\gamma}e(\eta).
\]
A presentation means the free unital $*$-algebra modulo the
$*$-ideal of these relations. Put $\mu_{p^r}=u^r$.
Its time evolution is
\begin{equation}
 \alpha_t(u)=p^{it}u,\qquad
 \alpha_t(e(\gamma))=e(\gamma),\qquad
 \alpha_t(u^*)=p^{-it}u^*.
 \label{v4-arith:kms-gns:eq:sigma-time}
\end{equation}
Every relation is homogeneous in occupation degree, so these
formulas define automorphisms. Finite algebra expressions have
an entire continuation in $t$.

For the fixed exponential character of
\Cref{v4-arith:thermal:prop:phase}, there is a surjective
$*$-homomorphism
\[
 \pi_+:\mathcal B_p\longrightarrow\mathcal T_p,
 \qquad u\mapsto S_p,\quad e(\gamma)\mapsto D_\gamma.
\]
Indeed every displayed relation holds on the occupation basis,
and $S_p,R_p$ generate the target. Define
$\varphi_{p,\beta}=\Phi_{p^{-\beta}}\circ\pi_+$.
This is a positive normalized functional: it is the pullback of
the explicitly summed positive Gibbs functional. It also extends
to the operator-norm completion of this represented quotient.
No faithfulness of $\pi_+$ is assumed.

\begin{lemma}[local KMS values]
\label{v4-arith:kms-gns:lem:KMS-property}
\ClaimStatusProvedHere. For real $\beta>0$, the functional satisfies
\begin{equation}
 \varphi_{p,\beta}(xy)
   =\varphi_{p,\beta}(y\alpha_{i\beta}(x))
       \qquad(x,y\in\mathcal B_p).
 \label{v4-arith:kms-gns:eq:KMS}
\end{equation}
For $n,m\in p^{\mathbb N_0}$ and $p^h\gamma=0$, putting
$q=p^{-\beta}$ gives
\begin{align}
 \varphi_{p,\beta}(\mu_n^*\mu_m e(\gamma))
   &=\delta_{nm}\varphi_{p,\beta}(e(\gamma)),
       \label{v4-arith:kms-gns:eq:phi-1}\\
 \varphi_{p,\beta}(e(\gamma))
   &=(1-q)\sum_{j=0}^{h-1}q^j e^{2\pi i p^j\gamma}+q^h.
       \label{v4-arith:kms-gns:eq:phase-state}
\end{align}
In particular $\varphi_{p,\beta}(\mu_m^*\mu_m)=1$.
\end{lemma}

\begin{proof}
The representation intertwines the time evolutions, so the
KMS identity follows from the normal-word calculation in
\Cref{v4-arith:thermal:prop:thermal-values}. If $n\ne m$,
the represented product has nonzero occupation degree and zero
diagonal entries. If $n=m$, its isometry product is one.
The finite phase expansion and its geometric tail give the
second formula. Taking $\gamma=0$ gives normalization at the
identity.
\end{proof}

Even uniqueness of a local state must retain its algebra.
Replacing the exponential character by its conjugate defines
another representation $\pi_-$ and another normalized KMS
functional. At $p=3$, $\beta=1$,
\[
 \varphi_{3,1}^+(e(1/3))=i\sqrt3/3,
 \qquad \varphi_{3,1}^-(e(1/3))=-i\sqrt3/3.
\]
Both values follow from the same finite phase sum. Thus the
presented local algebra has at least these two KMS states at
$\beta=1$. Also $\pi_+$ is not faithful at $p=3$: the element
\[
 e(1/3)-1-(e^{2\pi i/3}-1)(1-uu^*)
\]
vanishes under $\pi_+$ but not under $\pi_-$, as evaluation on
the vacuum shows. A statement about a particular represented
prime quotient cannot be substituted for a statement about the
full arithmetic algebra.

\subsection{GNS carrier}
\label{v4-arith:kms-gns:cycle1:GNS}

Use the inner product linear in the second variable,
$\langle[x],[y]\rangle=\varphi_{p,\beta}(x^*y)$.
The GNS null space is exactly $\ker\pi_+$, because the Gibbs
state on the represented algebra is faithful by
\eqref{v4-arith:thermal:gns-norm}. The quotient therefore has
precisely the GNS completion already constructed for $\mathcal T_p$.

\begin{lemma}[the represented local GNS space]
\label{v4-arith:kms-gns:lem:GNS-space}
\ClaimStatusProvedHere. Let $\varrho_q=(1-q)q^{N_p}$ and
$d_j=(1-q)q^j$. The unitary map and inner products are
\begin{equation}
 \mathscr G_{p,\beta}\cong\mathcal S_2(\mathscr H_p),
 \qquad [x]\longmapsto\pi_+(x)\varrho_q^{1/2},
 \qquad \Omega_{p,\beta}\longmapsto\varrho_q^{1/2},
 \label{v4-arith:kms-gns:eq:GNS-iso}
\end{equation}
\begin{equation}
 \langle[E_{ij}],[E_{kl}]\rangle
      =\delta_{ik}\delta_{jl}d_j,
 \qquad \langle[\mu_n],[\mu_m]\rangle=\delta_{nm}.
 \label{v4-arith:kms-gns:eq:GNS-inner-product}
\end{equation}
Here matrix units denote their classes in the represented
quotient. The local strong closure is left multiplication by
$B(\mathscr H_p)$, and the vacuum is cyclic and separating.
\end{lemma}

\begin{proof}
The null-space observation identifies the dense pre-Hilbert
quotient with the one in \Cref{v4-arith:thermal:gns-carrier}.
That theorem proves the unitary map, its dense range, all the
inner products and the strong closure. In particular the creation
vectors have norm one. The weights $m^{-\beta}$ occur on the
division vectors $[\mu_m^*]$, whose inner products are
$\delta_{nm}m^{-\beta}$. No independent orthogonal phase-coordinate
factor is obtained: phase inner products are the finite sums
\eqref{v4-arith:kms-gns:eq:phase-state}.
\end{proof}

\subsection{Modular sign and half-density}
\label{v4-arith:kms-gns:cycle1:modular}

\begin{lemma}[modular operator on the actual GNS carrier]
\label{v4-arith:kms-gns:lem:Delta}
\ClaimStatusProvedHere. On Hilbert--Schmidt matrix units,
\begin{equation}
 \Delta_{p,\beta}E_{ij}=p^{-\beta(i-j)}E_{ij},
 \label{v4-arith:kms-gns:eq:Delta-basis}
\end{equation}
with the maximal domain of \Cref{v4-arith:thermal:modular}.
Consequently
\begin{equation}
 \Delta_{p,\beta}^{it}[x]=[\alpha_{-\beta t}(x)],
 \qquad \Delta_{p,\beta}[\mu_m]=m^{-\beta}[\mu_m].
 \label{v4-arith:kms-gns:eq:Delta-formula}
\end{equation}
\end{lemma}

\begin{proof}
The ratio of the density weights in rows $i$ and $j$ is
$d_i/d_j=p^{-\beta(i-j)}$. The closed diagonal construction and
its domain are proved in \Cref{v4-arith:thermal:modular}. Its
imaginary powers therefore scale a degree-$r$ observable by
$p^{-it\beta r}$, establishing the sign in the time parameter.
\end{proof}

The densely defined conjugate-linear map and its closure are
\begin{equation}
 \mathfrak S_0[x]=[x^*],\qquad
 \overline{\mathfrak S_0}=J_{p,\beta}\Delta_{p,\beta}^{1/2}.
 \label{v4-arith:kms-gns:eq:S-def}
\end{equation}
The closure and polar decomposition, including the weighted
square-summability domain, were constructed by finite matrix
truncations in \Cref{v4-arith:thermal:modular}.

\begin{proposition}[modular conjugation of prime observables]
\label{v4-arith:kms-gns:prop:J-equals-sigma-BC}
\ClaimStatusProvedHere. The antiunitary $J_{p,\beta}$ is adjoint
on the Hilbert--Schmidt carrier. On cyclic vectors,
\begin{align}
 J_{p,\beta}[\mu_m]&=m^{\beta/2}[\mu_m^*],
       &J_{p,\beta}[\mu_m^*]&=m^{-\beta/2}[\mu_m],
       \label{v4-arith:kms-gns:eq:J-basis}\\
 J_{p,\beta}[e(\gamma)]&=[e(-\gamma)].
       \label{v4-arith:kms-gns:eq:J-on-algebra}
\end{align}
\end{proposition}

\begin{proof}
Apply \eqref{v4-arith:kms-gns:eq:S-def} to the eigenvectors of
\eqref{v4-arith:kms-gns:eq:Delta-formula}. The positive factor on
$[\mu_m]$ is $m^{-\beta/2}$, hence its inverse is the factor in
$J_{p,\beta}[\mu_m]$. Phases have degree zero and are fixed by
the modular positive factor. This also follows from the explicit
adjoint formula on matrices.
\end{proof}

\begin{theorem}[local KMS/GNS normalization]
\label{v4-arith:kms-gns:thm:KMS1-GNS-BC-identification}
\ClaimStatusProvedHere. For the represented prime Gibbs state,
$J_{p,\beta}$ is the antiunitary factor of the closed map
$[x]\mapsto[x^*]$. It fixes the vacuum and sends the left strong
operator algebra onto its commutant. At $\beta=1$, the factor on
$\mu_m\Omega$ is $m^{1/2}$ and the factor on $\mu_m^*\Omega$
is $m^{-1/2}$. This theorem constructs the modular normalization
on the stated quotient and Hilbert space. A Koszul--Petersson
identification additionally requires maps from the native pairing
carriers which intertwine these operations.
\end{theorem}

\begin{proof}
The Hilbert--Schmidt construction proves all operator assertions,
including the commutant, in \Cref{v4-arith:thermal:gns-carrier}
and \Cref{v4-arith:thermal:modular}. The half-density is also
forced by $\|\mu_m\Omega\|=1$ and
$\|\mu_m^*\Omega\|=m^{-\beta/2}$. Thus the inverse coefficient
on a creation vector would violate antiunitarity. None of these
calculations supplies a map from an arithmetic convolution category
or a Whittaker model into this representation.
\end{proof}

\begin{remark}[the required pairing map]
\label{v4-arith:kms-gns:rem:uniqueness}
The weighted lattice with squared basis norms $m^{-1}$ embeds
isometrically by sending its basis vector at $m=p^r$ to
$[\mu_m^*]$ at $\beta=1$. Sending it to $[\mu_m]$ with the
same normalization is not isometric. The first embedding reverses
occupation energy. The second can be made isometric by multiplying
$[\mu_m]$ by $m^{-1/2}$, but then it is a different comparison map.
Their image of the involution must be tested separately. The
uniqueness of the polar factor on a specified GNS space cannot
choose a native arithmetic embedding.
\end{remark}

\subsection{Prime-set assembly}
\label{v4-arith:kms-gns:cycle1:adelic}

\begin{corollary}[tensor shift GNS conjugation]
\label{v4-arith:kms-gns:cor:adelic-KMS-GNS}
\ClaimStatusProvedHere. For finite prime sets, the normalized
Gibbs GNS spaces tensor, and their modular conjugations tensor
with respect to the cyclic unit vectors. The resulting isometric
prime-set maps define an antiunitary involution on the incomplete
Hilbert tensor product of these spaces. This is a construction
for the tensor shift algebras and their specified reference states.
It does not identify the genuine arithmetic phase algebra or the
native adelic Hochschild object with that tensor product.
\end{corollary}

\begin{proof}
The last construction in \Cref{v4-arith:thermal:gns-section}
uses $X\mapsto X\otimes\varrho_{\mathrm{new}}^{1/2}$.
The new factor has norm one and is fixed by adjoint, so the
conjugations agree on the isometric inclusions and extend by
continuity to the Hilbert completion. The two-prime calculation
of \Cref{v4-arith:thermal:prop:two-prime-phase} proves that the
represented local phase quotients do not admit the naive
phase-preserving inclusions. Moreover, the occupation partition
at all primes diverges at $\beta=1$. These are distinct obstacles
to the stronger adelic identification.
\end{proof}

\section{Finite-Place Pairings and Representation Domains}
\label{v4-arith:kms-gns:cycle2}
\label{v4-arith:kms-gns:cycle2:ramified}

The original finite-place target concerns a native representation
factor, its Koszul dual, and a Whittaker or Kirillov function model.
The radial shell comparison of
\Cref{v4-arith:kp:radial-pairing-realization} constructs those
pairing maps on one explicit function subspace. Extending it to
an unramified or ramified principal series, or to a supercuspidal
factor, requires a map on the actual representation domain.
The state normalization alone cannot provide such an extension.

\begin{lemma}[linearity of local integral comparisons]
\label{v4-arith:kms-gns:lem:GJ-ramified}
A nonzero local zeta integral $Z(s,W,\Phi)$ linear in the vector
$W$ cannot have a formula
\begin{equation}
 Z(s,W,\Phi)=L(s)\int|\widehat\Phi(a)|^2\,d\mu_s(a)
 \label{v4-arith:kms-gns:eq:GJ-ramified}
\end{equation}
for every $W$ with fixed $\Phi$, unless the right side vanishes.
If a pairing is positive Hermitian, multiplying it by a scalar
$c$ preserves that property only when $c$ is positive real.
\end{lemma}

\begin{proof}
Replacing $W$ by $2W$ doubles the left side and leaves the
right side unchanged. For the pairing claim, evaluate on any
vector of positive squared norm. Its rescaled squared norm is
$c$ times a positive real number, so positivity forces
$c>0$. Conversely a positive real multiple preserves positivity.
\end{proof}

Thus a local functional equation must retain its vector and test
function dependence. Its complex root number cannot simply be
inserted as a common multiplier of two positive Hermitian norms.
Scalar normalizations act differently on the two kinds of forms:
rescaling one input of a bilinear form by $c$ multiplies its value
by $c$, whereas rescaling a vector map in a Hermitian norm
multiplies its squared norm by $|c|^2$.

\begin{proposition}[local realization with specified normalizations]
\label{v4-arith:kms-gns:lem:Koszul-ramified}
\label{v4-arith:kms-gns:lem:Pet-ramified}
\label{v4-arith:kms-gns:prop:local-match-ramified}
\ClaimStatusProvedHereConditional. Let $V_{p,\pi}$ be a chosen
native local representation core, with its dual and pairing
maps satisfying \Cref{v4-arith:kp:hyp:native-pairing-maps}.
Then its pairings obey
\begin{align}
 B_{p,\pi}(v,a)&=B_q(f_{p,\pi}v,g_{p,\pi}a),
    \label{v4-arith:kms-gns:eq:Koszul-ramified}\\
 h_{p,\pi}(v,w)&=h_q(f_{p,\pi}v,f_{p,\pi}w)
                =B_{p,\pi}(v,c_{p,\pi}w).
    \label{v4-arith:kms-gns:eq:Pet-ramified}
\end{align}
This statement applies to a ramified principal-series core when
those maps, including its local measure and scalar normalizations,
are constructed. It does not assert that they exist for every
ramified factor.
\end{proposition}

\begin{proof}
The first equality and the Hermitian equality are the specified
realization conditions. Their compatibility with the native
conjugate-linear map is the proved identity
$B_q(X,J_qY)=h_q(X,Y)$, giving the last equality. This is
\Cref{v4-arith:thermal:pairing-transfer}, with every carrier
and scalar normalization fixed before composition.
\end{proof}

\subsection{Supercuspidal and non-tempered domains}
\label{v4-arith:kms-gns:cycle2:supercuspidal}
\label{v4-arith:kms-gns:cycle2:ramanujan}

\begin{corollary}[supercuspidal conditional transfer]
\label{v4-arith:kms-gns:lem:supercuspidal}
\ClaimStatusProvedHereConditional. If the native pairing maps
are supplied on a chosen supercuspidal factor, the same local
pairing identity holds on that factor.
\end{corollary}
\begin{proof}
The proof of the local transfer uses the displayed maps and
pairings, and no principal-series presentation. It therefore
applies to that specified core.
\end{proof}

\begin{proposition}[the role of temperedness in a transfer]
\label{v4-arith:kms-gns:prop:non-tempered-residue}
\ClaimStatusProvedHereConditional. Suppose compatible native
pairing maps and their adelic assembly are constructed for every
factor in a specified tempered cuspidal domain. Any theorem
placing further automorphic representations in that domain allows
the same transfer on those representations. In particular a
Ramanujan-type temperedness hypothesis can enlarge that domain;
it cannot replace the native pairing-map construction.
\end{proposition}
\begin{proof}
Membership in the specified domain permits application of its
constructed maps. The conclusion is their local transfer followed
by their assumed adelic assembly. Membership alone contains no
map between the native and GNS carriers.
\end{proof}

\begin{remark}[modular spectrum and automorphic bounds]
\label{v4-arith:kms-gns:rem:partial-KS}
Temperedness does not make the local modular operator bounded:
$\Delta_qE_{0j}=q^{-j}E_{0j}$ is unbounded for every
$0<q<1$. The GNS construction itself works at every positive
real inverse temperature. Bounds on automorphic Satake parameters
control a different representation problem and do not alter this
matrix calculation.
\end{remark}

\begin{theorem}[conditional tempered cuspidal pairing transfer]
\label{v4-arith:kms-gns:thm:ramified-cuspidal-tempered-extension}
\ClaimStatusProvedHereConditional. Assume the native pairing
maps of \Cref{v4-arith:kp:hyp:native-pairing-maps} on the
chosen tempered cuspidal domain, including all its ramified
factors and its convergent or specified regularized adelic
assembly. Then
\begin{equation}
 \langle v,w\rangle^{\mathrm{Pet}}
    =\langle v,cw\rangle^{\mathrm{Koszul}}
 \label{v4-arith:kms-gns:eq:ramified-cuspidal-match}
\end{equation}
on that domain, with $c$ the native assembled conjugate-linear map.
\end{theorem}
\begin{proof}
Apply the proved local transfer on every factor. The assumed
assembly identifies the two products of local pairings with their
global pairings, so equality persists on elementary tensors and
on the specified completion or regularized domain. No independent
identification of genuine arithmetic phases with the tensor shift
algebra is used.
\end{proof}

\begin{remark}[native comparison boundary]
\label{v4-arith:kms-gns:rem:w3-tightening}
The explicit local state and the radial shell realization do not
construct the maps for all native tempered cuspidal factors.
That extension and the genuine arithmetic phase assembly remain
part of the original Koszul--Petersson problem. The operator
normalization provides a test for such maps, including their
variance, norms and energy domains.
\end{remark}

\section{Langlands Decomposition of \(L^{2}([GL_{2}])\) and Placement of \(\mathcal{V}^{\mathrm{prim}}\)}
\label{v4-arith:kms-gns:cycle3}

\subsection{The Spectral Decomposition}
\label{v4-arith:kms-gns:cycle3:spectral}

The Langlands decomposition of \(L^{2}([GL_{2}])\), following
Langlands 1966 and Moeglin--Waldspurger 1995, is:

\begin{theorem}[spectral decomposition of \(L^{2}(\lbrack GL_{2}\rbrack)\)]
\label{v4-arith:kms-gns:thm:spectral-decomposition}
\ClaimStatusProvedElsewhere. There is an orthogonal decomposition of
\(G(\mathbb{A})\)-modules:
\begin{equation}
\label{v4-arith:kms-gns:eq:L2-decomposition}
L^{2}([GL_{2}])\;=\;L^{2}_{\mathrm{cusp}}([GL_{2}])\;\oplus\;L^{2}_{\mathrm{res}}([GL_{2}])\;\oplus\;L^{2}_{\mathrm{cont}}([GL_{2}]),
\end{equation}
where:
\begin{itemize}
\item \(L^{2}_{\mathrm{cusp}}([GL_{2}])\) is the cuspidal spectrum,
a direct sum of irreducible unitary representations
\(\pi=\bigotimes'_{v}\pi_{v}\) with \(\pi_{p}\) generic tempered at
almost all primes;
\item \(L^{2}_{\mathrm{res}}([GL_{2}])\) is the residual spectrum
(in the \(GL_{2}\) case generated by one-dimensional characters
pulled up from the determinant), consisting of \(GL_{2}(\mathbb{A})\)-reducible
automorphic representations from residues of Eisenstein series at
specific points;
\item \(L^{2}_{\mathrm{cont}}([GL_{2}])\) is the continuous spectrum,
realised as the direct integral
\(\int^{\oplus}_{\mathrm{Re}(s)=1/2}\mathrm{Ind}_{B}^{G}(\chi_{s})\,|\xi(2s)|^{-2}ds\)
for \(\chi_{s}=|\cdot|^{s}\otimes|\cdot|^{-s}\).
\end{itemize}
\end{theorem}

\begin{proof}
Langlands 1966 \emph{On the Functional Equations Satisfied by
Eisenstein Series} \S6--7 for the construction; Moeglin--Waldspurger
1995 \emph{Spectral Decomposition and Eisenstein Series for \(GL_{n}\)}
for the published \(GL_{n}\) exposition; for \(GL_{2}\) specifically see
Gelbart 1975 \emph{Automorphic Forms on Adele Groups} Ch.~II.
\end{proof}

\subsection{Placement of \(\mathcal{V}^{\mathrm{prim}}\)}
\label{v4-arith:kms-gns:cycle3:placement}

\begin{proposition}[a native spectral placement requires its maps]
\label{v4-arith:kms-gns:prop:Vprim-placement}
\ClaimStatusProvedHereConditional. Suppose a native discrete core
has a direct-sum decomposition $V_{\mathrm{cusp}}\oplus V_{\mathrm{res}}$
and injective maps $\iota_{\mathrm{cusp}},\iota_{\mathrm{res}}$
into the corresponding orthogonal $L^2$ subspaces. Then
\begin{equation}
 \iota=\iota_{\mathrm{cusp}}\oplus\iota_{\mathrm{res}}
   :V_{\mathrm{cusp}}\oplus V_{\mathrm{res}}
       \longrightarrow L^2_{\mathrm{cusp}}\oplus L^2_{\mathrm{res}}
 \label{v4-arith:kms-gns:eq:placement}
\end{equation}
is injective. A map from an Eisenstein parameter to distributions
is a separate map and has no $L^2$ codomain unless square
integrability is established. The local GNS construction supplies
none of these native spectral maps.
\end{proposition}
\begin{proof}
Orthogonality of the two target subspaces makes the kernel of
$\iota$ the direct sum of the two kernels, hence zero. The
distributional and Hilbert codomains differ by definition;
square integrability must be checked before a distributional
expression represents a Hilbert vector. The explicit GNS target
is instead a space of operators with left multiplication.
\end{proof}

\section{Eisenstein Continuum: Koszul--Petersson Discrepancy}
\label{v4-arith:kms-gns:cycle4}

\subsection{The Maass--Selberg Inner Product}
\label{v4-arith:kms-gns:cycle4:maass-selberg}

The Maass--Selberg relations compute the Petersson-style inner
product of two Eisenstein series. Following Langlands 1966 \S7 and
Arthur 1980 Duke Math.~J.~45 \S3:

\begin{lemma}[Maass--Selberg inner product]
\label{v4-arith:kms-gns:lem:maass-selberg}
\ClaimStatusProvedElsewhere. For Eisenstein series
\(E(g,s_{1},\phi_{1})\) and \(E(g,s_{2},\phi_{2})\) with
\(\phi_{1},\phi_{2}\) in the induced representation
\(\mathrm{Ind}_{B}^{G}(\chi)\) and \(\mathrm{Re}(s_{1})=\mathrm{Re}(s_{2})=1/2\),
their truncated inner product is
\begin{equation}
\label{v4-arith:kms-gns:eq:maass-selberg}
\bigl\langle E(g,s_{1},\phi_{1}),E(g,s_{2},\phi_{2})\bigr\rangle^{\mathrm{Pet},\mathrm{trunc}}\;=\;\frac{1}{s_{1}+\overline{s_{2}}-1}\cdot\bigl\{\bigl\langle\phi_{1},\phi_{2}\bigr\rangle+\bigl\langle M(s_{1})\phi_{1},M(s_{2})\phi_{2}\bigr\rangle\bigr\},
\end{equation}
where \(M(s)\colon\mathrm{Ind}_{B}^{G}(\chi_{s})\to\mathrm{Ind}_{B}^{G}(\chi_{s}^{w})\)
is the Eisenstein intertwining operator with (standard) normalisation
\(M(s)=\xi(2s-1)/\xi(2s)\) in the unramified setting (Moeglin--Waldspurger
\S II.1.7).
\end{lemma}

\begin{proof}
Langlands 1966 \S7 (initial computation); Moeglin--Waldspurger 1995
\S II.1.7 (modern exposition); Arthur 1980 \S3 (adelic assembly).
The diagonal limit at \(s_{1}=s_{2}=s\), after removing the universal
truncation divergence, gives the Maass--Selberg scattering
normalisation at the Eisenstein parameter \(s\); it is not an ordinary
\(L^{2}\) density before this regularisation.
\end{proof}

\subsection{Thermal ratios and native continuation}
\label{v4-arith:kms-gns:cycle4:koszul-continuation}

The local thermal calculation supplies a scalar ratio of traces. It
must first be distinguished from a pairing on induced representations.
For \(\operatorname{Re}u>1\), write
\(Z_p(u)=\sum_{r\ge0}p^{-ur}=(1-p^{-u})^{-1}\). This is an ordinary
operator trace; it is a state partition function only for real positive
inverse temperature.

\begin{lemma}[the scalar Euler ratio]
\label{v4-arith:kms-gns:lem:Koszul-continued}
For \(\operatorname{Re}u>2\), both global occupation traces converge and
\[
 \prod_p\frac{Z_p(u)}{Z_p(u-1)}
 =\prod_p\frac{1-p^{1-u}}{1-p^{-u}}
 =\frac{\zeta(u)}{\zeta(u-1)}.
\]
Define \(\Lambda(u)=\pi^{-u/2}\Gamma(u/2)\zeta(u)\). Its scalar
completed ratio is
\begin{equation}
 R(u)=\frac{\Lambda(u)}{\Lambda(u-1)}
 =\pi^{-1/2}\frac{\Gamma(u/2)}{\Gamma((u-1)/2)}
                 \frac{\zeta(u)}{\zeta(u-1)}.
 \label{v4-arith:kms-gns:eq:Koszul-continued}
\end{equation}
The meromorphic function \(R\) has a removable singularity at \(u=1\)
with value \(-1\). The ratio of the entire completed functions
\(\xi_{\mathbb R}(u)/\xi_{\mathbb R}(u-1)\) instead has value \(1\)
there. Neither has the asserted simple pole at that point.
\end{lemma}
\begin{proof}
Absolute convergence for \(\operatorname{Re}u>2\) follows by summing
\(p^{1-\operatorname{Re}u}\) and \(p^{-\operatorname{Re}u}\).
Dividing the two absolutely convergent Euler products gives the first
identity. Multiplication by the gamma factors gives the second.
The zeta residue at one is one, so
\(\Lambda(1+z)=z^{-1}+O(1)\). The functional equation
\(\Lambda(u)=\Lambda(1-u)\) gives
\(\Lambda(z)=-z^{-1}+O(1)\), hence \(R(1+z)\to-1\).
Since \(\xi_{\mathbb R}(u)=\tfrac12u(u-1)\Lambda(u)\), its values at
zero and one are both \(1/2\), proving the last assertion.
\end{proof}

Analytic continuation extends these scalar functions. It does not
extend a Gibbs density through the divergent global partition function
at \(\beta=1\), and it does not construct a Koszul pairing on an induced
representation. The latter requires maps from that representation and
its native dual into specified pairing carriers. For a Hermitian
parameter family the complex variables are \((s_1,\overline{s_2})\);
meromorphy in those variables is not meromorphy in \((s_1,s_2)\).

\begin{theorem}[conditional comparison of scalar normalizations]
\label{v4-arith:kms-gns:thm:eisenstein-discrepancy}
Suppose on a specified pair of parameter-dependent lines a common
sesquilinear form \(h_0\), a regularized automorphic pairing \(h_P\), and
a native bilinear pairing \(B\) with conjugate-linear duality \(c\) are
constructed. Suppose, on their common regular domain, their scalar
normalizations are
\(h_P=P(s_1,\overline{s_2})h_0\) and
\(B(\cdot,c\,\cdot)=R(s_1+\overline{s_2})h_0\). Then
\begin{equation}
 h_P-B(\cdot,c\,\cdot)
 =\bigl(P(s_1,\overline{s_2})-R(s_1+\overline{s_2})\bigr)h_0.
 \label{v4-arith:kms-gns:eq:continuum-discrepancy}
\end{equation}
At a pair where \(h_0\ne0\), equality of pairings is equivalent to
\(P(s_1,\overline{s_2})=R(s_1+\overline{s_2})\).
In particular, if \(P=A/(u-1)\), the difference is
\((A/(u-1)-R(u))h_0\), not \((A-R(u))h_0/(u-1)\).
\end{theorem}
\begin{proof}
Subtract the two assumed scalar multiples. Multiplication by a
nonzero value of \(h_0\) is injective on the scalar field. The last
formula is the same subtraction with the specified denominator.
No pole cancellation or equality of residues follows unless the two
Laurent expansions are also supplied.
\end{proof}

\begin{proposition}[a deciding scalar test on the critical line]
\label{v4-arith:kms-gns:rem:match-is-FE}
Put \(M(s)=\Lambda(2s-1)/\Lambda(2s)\). At every point
\(s=1/2+it\) where both factors are finite and nonzero,
\(\lvert M(s)\rvert=\lvert R(2s)\rvert=1\). Consequently the proposed
scalar identity
\begin{equation}
       1+|M(s)|^2=R(2s)
       \label{v4-arith:kms-gns:eq:continuum-match-condition}
\end{equation}
cannot hold at such a point.
\end{proposition}
\begin{proof}
The real coefficients and completed functional equation give
\(\Lambda(2it)=\Lambda(1-2it)=\overline{\Lambda(1+2it)}\).
Thus \(M(s)\) is a quotient of a nonzero complex number's conjugate by
that number and has modulus one; \(R(2s)=M(s)^{-1}\).
The left side of the proposed identity is two and the right side has
modulus one. This refutes the scalar identity, independently of any
native pairing realization.
\end{proof}

\begin{corollary}[boundary of the continuum comparison]
\label{v4-arith:kms-gns:cor:eisenstein-final}
The thermal ratios do not determine a native Eisenstein pairing, and
the displayed scalar test does not decide equality of the actual native
and automorphic pairings. That equality requires the maps, domains,
normalizations and regularized spectral measures used in the
conditional comparison theorem.
\end{corollary}
\begin{proof}
The trace ratio is a scalar function on the occupation representation.
Its construction includes no map from an induced representation or its
Koszul dual. The conditional theorem uses those maps as hypotheses;
its proof does not produce them. The scalar contradiction therefore
excludes the proposed normalization identity, not an independently
constructed pairing with a different derived normalization.
\end{proof}

\begin{remark}[wave packets and their domains]
\label{v4-arith:kms-gns:rem:continuum-honest-domain}
Individual Eisenstein parameters generally represent distributions,
whereas their suitably integrated wave packets can represent Hilbert
vectors. To compare pairings on the continuum one must specify the test
functions, convergence or truncation, the two spectral distributions,
and their equality on those tests. A meromorphic scalar ratio alone
establishes none of these statements. The conditional discrete-spectrum
transfer does not use this continuum comparison.
\end{remark}

\section{Consequence for P3}
\label{v4-arith:kms-gns:P3}

\begin{theorem}[positive pairing transfer on a realized core]
\label{v4-arith:kms-gns:thm:P3-tempered-cuspidal}
\ClaimStatusProvedHereConditional. Suppose the chosen native
cuspidal core has an injective square-integrable realization and
satisfies the pairing-map and assembly hypotheses of
\Cref{v4-arith:kms-gns:thm:ramified-cuspidal-tempered-extension}.
Then its Petersson form is positive definite and equals
$\langle v,cw\rangle^{\mathrm{Koszul}}$ on that core. This
constructs a positive Hermitian form under the specified maps;
it asserts no vertex-algebra structure or full-sector equivalence.
\end{theorem}
\begin{proof}
For a nonzero native vector, injectivity gives a nonzero $L^2$
function, whose integral squared norm is positive. Equality of
pairings is the conditional transfer already proved. The local
GNS positivity does not remove the hypothesis of a native
square-integrable realization.
\end{proof}

\begin{remark}[domains of the pairing comparison]
\label{v4-arith:kms-gns:rem:P3-domain-enlargement}
The tempered cuspidal transfer requires native maps on that
core. Further automorphic temperedness results do not create
those maps. Eisenstein pairings require separate domains and
regularizations. Neither the local KMS theorem nor an abstract
Hilbert-space isomorphism supplies a full-sector extension.
\end{remark}

\subsection{The arithmetic implication}
\label{v4-arith:kms-gns:P3:rh}
\begin{corollary}[the pairing input remains conditional]
\label{v4-arith:kms-gns:cor:rh-minimal-updated}
The local GNS normalization does not close the native
Koszul--Petersson input in a Hilbert--P\'olya or positivity
argument. Such an implication retains the native pairing maps,
their actual domain and their assembly as hypotheses.
\end{corollary}
\begin{proof}
The local theorem constructs a represented quotient and its
operator Hilbert space. The native transfer theorem explicitly
requires the additional maps before its equations can be
composed. Omitting them removes the first equality connecting
native and represented pairings.
\end{proof}

\section{The Remaining Native Construction}
\label{v4-arith:kms-gns:residual-list}
\label{v4-arith:kms-gns:verification}
\label{v4-arith:kms-gns:summary}

The prime-power occupation algebra, its normalized state,
Hilbert--Schmidt GNS realization, modular domains and conjugation
are explicit. The radial shell model realizes a bilinear pairing,
a Hermitian pairing and their anti-linear comparison by two
separate maps with opposite dual energy degrees. These give
concrete local constructions and decisive tests of normalization.

The original native problem still requires maps from the
arithmetic Hochschild and Koszul factors to the chosen operator
and Whittaker carriers, with their actual Hecke and phase actions.
It also requires a compatible adelic assembly. The represented
single-prime phase quotients cannot supply that assembly by the
naive inclusions, and a global occupation density at $\beta=1$
is not nuclear. The finite-prime tensor GNS construction is
therefore not an adelic arithmetic equivalence.

The automorphic and continuous-spectrum questions remain
separate. A local functional equation has to retain its vectors,
test functions and complex normalization factors. A Hermitian
norm has positive real normalization. A spectral distribution
needs its test space and regularization before it can be compared
with a Hilbert pairing. The GNS calculation neither proves these
native comparisons nor decides positivity for a distinct
arithmetic Hilbert--P\'olya operator.
